\section{Data and exploratory analysis}
\label{sec:data}

\subsection{Corpora}
All corpora are downloaded inside the notebook from public Hugging Face identifiers (Table~\ref{tab:corpora}). The three label structures differ by design: topic, polarity and intent.

\begin{table}[H]
\centering
\footnotesize
\caption{Corpora. AG News is stratified-sampled from 120{,}000 training and 7{,}600 test texts; the others use the official splits.}
\label{tab:corpora}
\begin{tabular}{@{}llrP{4.6cm}@{}}
\toprule
Corpus & Hugging Face identifier & Classes & Label meaning \\
\midrule
AG News & {\scriptsize\texttt{fancyzhx/ag\_news}} & 4 & World, Sports, Business, Sci/Tech \\
Rotten Tomatoes & {\scriptsize\texttt{cornell-movie-review-data/rotten\_tomatoes}} & 2 & negative / positive review \\
Banking77 & {\scriptsize\texttt{mteb/banking77}} & 77 & customer-service intent \\
\bottomrule
\end{tabular}
\end{table}

\subsection{Data quality and cleaning}
No corpus contains missing or empty texts. AG News stores HTML entities without the leading ampersand (\texttt{Google \#39;s}, \texttt{quot;}, \texttt{\&lt;b\&gt;}) and marks line breaks with a backslash; before cleaning, the entity number \texttt{39} was among the ten most frequent tokens. The cleaning step restores these characters, removes markup and normalises whitespace for all corpora. The AG News training split contains 156 exact duplicates; Banking77 has 7 development texts that also occur in the test split, which are removed from fit and validation (Table~\ref{tab:eda}).

\begin{table}[H]
\centering
\footnotesize
\caption{Fit-split statistics after cleaning (lengths in regex word tokens; class size = min--max examples per class; Neg.\,\% = texts with a negation cue; Dup. = duplicate texts in the original training split; Leak rm. = development texts removed because they also occur in a later split).}\label{tab:eda}
\begin{tabular}{@{}lrrrrrrrrr@{}}
\toprule
Corpus & Fit / Val / Test & $K$ & Class size & Mean len & P95 & Vocab & Neg.\,\% & Dup. & Leak rm. \\
\midrule
AG News & 8,500 / 1,500 / 2,000 & 4 & 2125--2125 & 38.2 & 54 & 23,901 & 13.7 & 156 & 0 \\
Rotten Tomatoes & 7,250 / 1,280 / 1,066 & 2 & 3625--3625 & 18.9 & 34 & 15,759 & 23.4 & 0 & 0 \\
Banking77 & 8,488 / 1,498 / 3,076 & 77 & 30--159 & 12.0 & 29 & 2,243 & 25.9 & 0 & 7 \\
\bottomrule
\end{tabular}
\end{table}


\subsection{Observations}
Figures~\ref{fig:eda-ag}--\ref{fig:eda-b77} show, for each fit split, the sorted class sizes, the length distribution and the most frequent words after stop-word removal.

\begin{figure}[H]
\centering
\includegraphics[width=\textwidth]{ag_news/eda.png}
\caption{AG News (fit split). Balanced classes; news-agency and date words dominate the content vocabulary.}
\label{fig:eda-ag}
\end{figure}

\begin{figure}[H]
\centering
\includegraphics[width=\textwidth]{rotten_tomatoes/eda.png}
\caption{Rotten Tomatoes (fit split). Balanced polarity; frequent content words (\textit{film, movie, like, good}) are not tied to one polarity.}
\label{fig:eda-rt}
\end{figure}

\begin{figure}[H]
\centering
\includegraphics[width=\textwidth]{banking77/eda.png}
\caption{Banking77 (fit split). Class sizes vary five-fold (30--159); a few domain words (\textit{card, account, money, transfer}) are shared by many intents.}
\label{fig:eda-b77}
\end{figure}

\begin{itemize}
  \item \textbf{Length.} Texts are short: a mean of 38, 19 and 12 words, with 95th percentiles of 54, 34 and 29 words. A token budget of 96/64 therefore truncates almost nothing, and the sequence models see nearly every word.
  \item \textbf{Balance.} AG News and Rotten Tomatoes are perfectly balanced, so accuracy and macro-F1 almost coincide. Banking77 is moderately imbalanced in training (ratio 5.3) but balanced in test (about 40 texts per intent); macro-F1 is the appropriate summary.
  \item \textbf{Vocabulary.} AG News has the richest vocabulary (23{,}901 types) and topic-specific nouns (\textit{oil}, team names, companies). Banking77 has only 2{,}243 types: intents are distinguished less by \emph{which} words occur than by \emph{how} they combine (\textit{card arrival} vs.\ \textit{card delivery estimate}; \textit{top up failed} vs.\ \textit{top up reverted}).
  \item \textbf{Function words carry meaning.} Negation cues occur in 13.7\%, 23.4\% and 25.9\% of the texts. In Banking77 the most frequent words are \textit{i, my, to, can, do}; in Rotten Tomatoes polarity often hinges on \textit{not, but, too}. The standard English stop-word list removes all of these (it contains \textit{not, no, never, nor, cannot, how, why, what}).
\end{itemize}

\paragraph{Expectations derived from the EDA.}
\textbf{E1:} on AG News, topical nouns should make TF-IDF competitive with pretrained embeddings.
\textbf{E2:} on Rotten Tomatoes and Banking77, stop-word removal should hurt TF-IDF, and bigrams should help.
\textbf{E3:} pretrained sentence embeddings should help most where labels depend on composition or paraphrase (Rotten Tomatoes, Banking77).
\textbf{E4:} a model trained from scratch on ${\approx}$8{,}500 texts (about 110 per class in Banking77, 3{,}600 in Rotten Tomatoes) should trail both feature-based pipelines.
